\documentclass[aspectratio=169,11pt]{beamer}

\usetheme{Madrid}
\usecolortheme{default}
\usefonttheme{professionalfonts}
\setbeamertemplate{navigation symbols}{}
\setbeamertemplate{footline}[frame number]

\usepackage{amsmath, amssymb, booktabs, graphicx}

\newcommand{\E}{\mathbb{E}}
\newcommand{\1}{\mathbf{1}}

\title{Prediction, Regularization, and Applied Measurement}
\subtitle{Empirical Methods --- Week 9}
\author{Teaching deck draft}
\date{}

\begin{document}

\begin{frame}
  \titlepage
\end{frame}

\begin{frame}{Week 9 question}
\begin{itemize}
    \item When is prediction itself the research object?
    \item When is prediction an input into measurement or design?
    \item How do train/test logic, regularization, and calibration discipline the object?
    \item What can go wrong when a predicted variable enters an economics paper?
\end{itemize}

\medskip
Machine learning can improve measurement and diagnostics, but it does not define the estimand or identify a causal effect.
\end{frame}

\begin{frame}{Prediction versus causal design}
\small
\begin{tabular}{p{0.28\linewidth} p{0.30\linewidth} p{0.30\linewidth}}
\toprule
Use & Main object & Evaluation question \\
\midrule
Prediction as final object & risk, forecast, ranking, probability & Does it generalize and calibrate on the target population? \\
Prediction as measurement & score, label, exposure, proxy & Does it measure the economic construct needed downstream? \\
Prediction as nuisance input & flexible controls or nuisance functions & Does it help estimate the causal object without becoming the answer? \\
\bottomrule
\end{tabular}

\medskip
\begin{itemize}
    \item Prediction accuracy is not identification.
    \item A good score can still be a bad measure if the target, timing, or population is wrong.
    \item The downstream design determines which prediction errors matter.
\end{itemize}
\end{frame}

\begin{frame}{Train/test risk}
\small
\[
R(f)=\E\left[L(Y,f(X))\right]
\]

\[
\widehat f
=
\arg\min_{f\in\mathcal F}
\frac{1}{|\mathcal I_{train}|}
\sum_{i\in\mathcal I_{train}} L(Y_i,f(X_i))
\]

\[
\widehat R_{test}(\widehat f)
=
\frac{1}{|\mathcal I_{test}|}
\sum_{i\in\mathcal I_{test}} L(Y_i,\widehat f(X_i))
\]

\begin{itemize}
    \item Random splits answer a narrow generalization question.
    \item Time, firm, region, or occupation splits test transportability.
    \item The loss should match the economics cost of the mistake.
\end{itemize}
\end{frame}

\begin{frame}{Regularization and tuning}
\small
\[
\widehat\beta^{lasso}
=
\arg\min_{\beta}
\frac{1}{n}\sum_{i=1}^{n}(Y_i-X_i'\beta)^2
+
\lambda\sum_{j=1}^{p}|\beta_j|
\]

\[
\widehat\beta^{ridge}
=
\arg\min_{\beta}
\frac{1}{n}\sum_{i=1}^{n}(Y_i-X_i'\beta)^2
+
\lambda\sum_{j=1}^{p}\beta_j^2
\]

\[
\widehat\beta^{EN}
=
\arg\min_{\beta}
\frac{1}{n}\sum_{i=1}^{n}(Y_i-X_i'\beta)^2
+
\lambda\left[\alpha\sum_j|\beta_j|+(1-\alpha)\sum_j\beta_j^2\right]
\]
\end{frame}

\begin{frame}{Cross-validation selects prediction rules}
\small
\[
\widehat\lambda
=
\arg\min_{\lambda\in\Lambda}
\frac{1}{K}\sum_{k=1}^{K}
\frac{1}{|V_k|}
\sum_{i\in V_k}
L\left(Y_i,\widehat f_{-k,\lambda}(X_i)\right)
\]

\medskip
\begin{tabular}{p{0.22\linewidth} p{0.32\linewidth} p{0.34\linewidth}}
\toprule
Penalty & Best when & Caveat \\
\midrule
Lasso & sparse signal is plausible & selected variables can be unstable \\
Ridge & many related predictors matter & no sparse interpretation \\
Elastic net & grouped or collinear features & tuning has two moving parts \\
\bottomrule
\end{tabular}

\medskip
Cross-validation is a tuning device. It is not an exclusion restriction, a research design, or a policy-invariance argument.
\end{frame}

\begin{frame}{Bias, variance, and calibration}
\small
\[
\E\left[(Y-\widehat f(x))^2\mid X=x\right]
=
\sigma^2+
\left(\E[\widehat f(x)]-f_0(x)\right)^2+
\operatorname{Var}(\widehat f(x))
\]

\[
\E[Y_i\mid \widehat p_i=s]=s,
\qquad
\widehat B=\frac{1}{n}\sum_{i=1}^{n}(Y_i-\widehat p_i)^2
\]

\[
\widehat C_i(t)=\1\{\widehat p_i\geq t\}
\]

\begin{itemize}
    \item Regularization accepts bias to reduce variance.
    \item Calibration matters when a score is used as probability or intensity.
    \item Thresholds encode policy or measurement tradeoffs.
\end{itemize}
\end{frame}

\begin{frame}{Economics applications to measurement}
\small
\begin{tabular}{p{0.24\linewidth} p{0.31\linewidth} p{0.33\linewidth}}
\toprule
Setting & Predicted object & Downstream use \\
\midrule
Job postings & skills, tasks, technology labels & labor demand, task change, technology exposure \\
Raw job titles & occupation or task group & mobility, wage growth, exposure by occupation \\
Remote work & feasibility or intensity score & spatial sorting, firm policy, worker access \\
Risk environments & default, attrition, take-up & targeting, diagnostics, policy triage \\
Latent outcomes & imputed score or proxy & measurement when outcomes are missing or costly \\
\bottomrule
\end{tabular}

\medskip
The question is not ``which algorithm is strongest?'' It is whether the predicted object matches the economic construct.
\end{frame}

\begin{frame}{Feature engineering and leakage}
\small
\begin{itemize}
    \item Predictors must be available at the time and level of the intended prediction.
    \item Leakage can enter through post-outcome variables, post-treatment text, edited metadata, future codes, or labels embedded in features.
    \item Job-posting text, occupation codes, firm histories, and worker records all require a timing diagram.
    \item A model can pass a random holdout test while failing the actual research-use test.
\end{itemize}

\medskip
\[
\text{allowed feature}
\quad \Longleftrightarrow \quad
X_i \text{ observed before } Y_i \text{ and before the policy/design decision.}
\]
\end{frame}

\begin{frame}{Implementation pitfalls}
\small
\begin{tabular}{p{0.25\linewidth} p{0.33\linewidth} p{0.30\linewidth}}
\toprule
Pitfall & Why it matters & Check \\
\midrule
Class imbalance & accuracy hides rare-event failure & precision, recall, confusion matrix \\
Drift & language and behavior change & time-split validation \\
Transportability & model fails across firms or regions & group-split validation \\
Instability & selected features are sample-specific & fold-level selection stability \\
Subgroup error & measurement error is systematic & subgroup loss and calibration \\
Interpretability & measure is not auditable & target, features, labels, failure modes \\
\bottomrule
\end{tabular}
\end{frame}

\begin{frame}{Research Lab design}
\begin{columns}[T,onlytextwidth]
\begin{column}{0.32\linewidth}
\textbf{Reproduce}
\begin{itemize}
    \item synthetic job postings
    \item keyword baseline
    \item ridge, lasso, elastic net
    \item cross-validation and test loss
\end{itemize}
\end{column}
\begin{column}{0.32\linewidth}
\textbf{Diagnose}
\begin{itemize}
    \item calibration bins
    \item subgroup performance
    \item leakage audit
    \item feature stability
\end{itemize}
\end{column}
\begin{column}{0.32\linewidth}
\textbf{Transfer}
\begin{itemize}
    \item occupation-title data
    \item vocabulary coverage
    \item transportability loss
    \item downstream-use memo
\end{itemize}
\end{column}
\end{columns}
\end{frame}

\begin{frame}{Takeaways}
\begin{itemize}
    \item Start with the economic object, not the algorithm.
    \item Train/test discipline estimates generalization only for the population implied by the split.
    \item Regularization is useful because high-dimensional data are noisy and collinear.
    \item A predicted measure needs calibration, subgroup diagnostics, leakage checks, and transportability evidence.
    \item ML improves applied economics when it sharpens measurement while leaving identification explicit.
\end{itemize}
\end{frame}

\end{document}
